\documentclass[11pt,a4paper]{article}
\usepackage[T1]{fontenc}
\usepackage{lmodern}
\usepackage{amsmath,amssymb,amsthm}
\usepackage[margin=28mm]{geometry}
\usepackage{microtype}
\usepackage[hidelinks]{hyperref}
\newtheorem{theorem}{Theorem}
\newtheorem{lemma}[theorem]{Lemma}
\newcommand{\Prob}{\mathbb P}
\newcommand{\E}{\mathbb E}
\newcommand{\Z}{\mathbb Z}
\title{CO phase selection for a stirred catalytic surface model}
\author{Kaiwen Jin\\[3pt]
  \small\href{mailto:yui.ui.siki@gmail.com}{yui.ui.siki@gmail.com}}
\date{15 October 2026}
\hypersetup{
  pdftitle={CO phase selection for a stirred catalytic surface model},
  pdfauthor={Kaiwen Jin}}

\begin{document}
\maketitle

\begin{abstract}
We consider a catalytic surface model on \(\Z^2\) with CO adsorption,
diatomic oxygen adsorption, finite-rate reaction and finite-rate stirring.
If the CO arrival rate \(p\) exceeds the oxygen parameter \(q\), we prove
that every site is eventually occupied by CO almost surely, from every
iid initial law with positive CO density. This gives CO poisoning on
every finite set. The density estimate of Durrett and Swindle controls
the total time spent vacant. An additional reaction count then controls
all state changes, including stirring. Spatial independence of the
initial states and driving clocks determines the limiting phase.
\end{abstract}

\section{The model and the result}

In the catalytic oxidation model of Durrett and Swindle \cite{DS}, a
surface site is vacant, occupied by a CO molecule, or occupied by an
oxygen atom. Neighboring CO and oxygen react and leave two vacant sites.
The two configurations consisting entirely of CO or entirely of oxygen
are absorbing; either one stops the oxidation reaction.

Write \(\eta_t(x)\in\{0,1,2\}\), \(x\in\Z^2\), for the state at time \(t\),
with 0 vacant, 1 CO and 2 oxygen. Sites \(x,y\) are neighbors if
\(|x-y|_1=1\). The transition rates are
\[
\begin{array}{c|c|c}
\text{location}&\text{update}&\text{rate}\\ \hline
\text{site}&0\longrightarrow1&p\\
\text{unordered neighboring pair}&(0,0)\longrightarrow(2,2)&q/4\\
\text{unordered neighboring pair}&(1,2),(2,1)\longrightarrow(0,0)&r/4\\
\text{unordered neighboring pair}&(a,b)\longrightarrow(b,a)&\gamma .
\end{array}
\]
In the last row \(a,b\in\{0,1,2\}\); exchanging equal states has no effect.
All rates below are finite and satisfy
\[
 p>q>0,\qquad r>0,\qquad \gamma>0.
\]
Each update has its own independent Poisson clock. The initial states are
independent of these clocks. Reaction is at a finite rate, rather than
instantaneous.

Durrett and Swindle \cite[Theorem~1 and its proof, pp.~490--492]{DS}
showed that the vacant-site and neighboring mixed-species probabilities
vanish when \(p>q\), under translation-invariant initial laws. Thus every
fixed finite box becomes monochromatic with probability tending to one.
Selecting the CO phase requires a further argument. It is posed in
\cite[Section~5.4, Open Problem~5.4.6]{IPS} for more general initial
configurations. Here we prove selection for iid initial laws of positive
CO density.

\begin{theorem}\label{thm:selection}
For the process above, let the initial states be iid with
\(\Prob(\eta_0(0)=1)>0\). Almost surely, for every \(x\in\Z^2\),
\(\eta_t(x)=1\) for all sufficiently large \(t\). Consequently, for every
finite \(K\subset\Z^2\),
\begin{equation}\label{eq:target}
 \lim_{t\to\infty}
 \Prob\bigl(\eta_t(x)=1\text{ for every }x\in K\bigr)=1.
\end{equation}
\end{theorem}

The time after which a site stays in state 1 may depend on the site.
The theorem does not assert that the whole infinite lattice is absorbed
at one finite time. Its conclusion holds for each fixed admissible
parameter choice, without a uniform convergence rate.

The proof retains the signed-density estimate from \cite{DS} and uses
the CO balance once more to bound the total number of reactions.
These two bounds control the expected number of each kind of state change.
Once all sites fix, connectedness leaves only the two constant configurations, and
the iid driving data make the choice of color deterministic. The
signed density then excludes oxygen.

\section{There are finitely many state changes}

We use the graphical process driven by the independent clocks described
above. It is a measurable, nonanticipating function of the initial states
and clocks and commutes with spatial translations. If \(C\) is any
driving clock of rate \(c\), \(F\subset\Z^2\) is finite and \(h\) is a
bounded function of the states in \(F\), then
\begin{equation}\label{eq:counts}
 \E\sum_{s\in C\cap(0,T]}h(\eta_{s-}|_F)
       =c\int_0^T \E h(\eta_t|_F)\,dt .
\end{equation}
Appendix~\ref{sec:clocks} gives the construction and a direct proof of
this counting identity. It applies both to state increments and to the
indicator that a clock ring changes a given site.

Let \(a_i=\Prob(\eta_0(0)=i)\) and
\(u_i(t)=\Prob(\eta_t(0)=i)\), \(i=0,1,2\). To count oxygen additions
and CO losses at the origin, put
\begin{align*}
 O(t)&=\frac q4\sum_{y\sim0}
       \Prob(\eta_t(0)=0,\eta_t(y)=0),\\
 R(t)&=\frac r4\sum_{y\sim0}
       \Prob(\eta_t(0)=1,\eta_t(y)=2).
\end{align*}
Every sum here has four terms. Translation invariance, with \(y\) replaced
by \(-y\), gives
\begin{equation}\label{eq:orientations}
 \sum_{y\sim0}\Prob(\eta_t(0)=2,\eta_t(y)=1)
 =\sum_{y\sim0}\Prob(\eta_t(0)=1,\eta_t(y)=2).
\end{equation}
Stirring contributes zero to a one-site density, since its contribution
for state \(i\) is
\(\gamma\sum_{y\sim0}[u_i(t)-u_i(t)]\).
Each effective oxygen ring on an incident edge adds one oxygen atom at
the origin. Applying \eqref{eq:counts} to the state indicators therefore
gives the exact balances
\begin{align}
 u_1(T)-a_1&=\int_0^T[p u_0(t)-R(t)]\,dt,\label{eq:CO}\\
 u_2(T)-a_2&=\int_0^T[O(t)-R(t)]\,dt.\label{eq:O}
\end{align}
No independence of neighboring states at positive times is used.

Set \(m(t)=u_1(t)-u_2(t)\). Since \(O(t)\le q u_0(t)\), subtraction yields
\begin{equation}\label{eq:drift}
 m(T)-m(0)
   =\int_0^T[p u_0(t)-O(t)]\,dt
   \ge(p-q)\int_0^T u_0(t)\,dt\ge0.
\end{equation}
This is the density argument in \cite{DS}. The bound \(m(T)\le1\) implies
\begin{equation}\label{eq:I}
 I:=\int_0^\infty u_0(t)\,dt
        \le\frac{1-m(0)}{p-q}<\infty.
\end{equation}
The CO balance \eqref{eq:CO} also gives
\begin{equation}\label{eq:J}
 J:=\int_0^\infty R(t)\,dt
      \le a_1+pI<\infty,
\end{equation}
because \(u_1(T)\ge0\). These limits use monotone convergence of
nonnegative integrals.

\begin{lemma}\label{lem:finite}
Under the hypotheses of Theorem~\ref{thm:selection}, almost surely every
site changes state only finitely many times.
\end{lemma}

\begin{proof}
Let \(N_0(T)\) count state changes at the origin up to time \(T\).
The expected numbers due to CO arrivals, oxygen arrivals and reactions
are, respectively,
\[
 p\int_0^T u_0(t)\,dt\le pI,\qquad
 \int_0^T O(t)\,dt\le qI,\qquad
 2\int_0^T R(t)\,dt\le2J.
\]
The factor 2 in the reaction count comes from the two possible colors
at the origin in \eqref{eq:orientations}.

For stirring let
\(D_y(t)=\Prob(\eta_t(0)\ne\eta_t(y))\).
An unequal pair either contains a vacancy or consists of a CO and
an oxygen atom, so
\[
 D_y(t)\le2u_0(t)
       +\Prob(\eta_t(0)=1,\eta_t(y)=2)
       +\Prob(\eta_t(0)=2,\eta_t(y)=1).
\]
Thus \eqref{eq:orientations} bounds the expected number of effective
exchanges at the origin by
\[
 \gamma\int_0^T\sum_{y\sim0}D_y(t)\,dt
        \le8\gamma I+\frac{8\gamma}{r}J.
\]
Every change has been counted. Letting \(T\) increase to infinity,
\begin{equation}\label{eq:allchanges}
 \E N_0(\infty)
   \le(p+q+8\gamma)I+(2+8\gamma/r)J<\infty.
\end{equation}
Hence \(N_0(\infty)\) is finite almost surely. Translation invariance
gives the same conclusion at each site; the countable union of the
exceptional null events still has probability zero.
\end{proof}

\section{Selection of the limiting color}

\begin{lemma}\label{lem:color}
Almost surely the eventual states of all sites are the same and belong
to \(\{1,2\}\).
\end{lemma}

\begin{proof}
By Lemma~\ref{lem:finite}, every site has an eventual state.
Simultaneously for all the countably many positive-rate clocks, their
ring times are unbounded almost surely. Indeed, for a rate-\(c\) clock
the probability of no ring after an integer \(n\) is
\(\lim_{M\to\infty}e^{-c(M-n)}=0\).

An eventual vacancy at a site would be changed by a later CO ring.
Two neighboring eventual states 1 and 2 would be changed by a reaction
ring after both sites had fixed. Neither is possible. Connectedness of
\(\Z^2\) now makes the eventual configuration constant. These are
sample-path arguments and require no stopping-time property of the
last-change times.
\end{proof}

\begin{proof}[Proof of Theorem~\ref{thm:selection}]
Let \(e_1,e_2\) be the coordinate unit vectors. At each \(x\), group its
initial state and site clock with all three clocks on each of
\(\{x,x+e_1\}\) and \(\{x,x+e_2\}\). Denote the resulting mark by \(Y_x\).
Every unordered edge is assigned once, so the marks
\((Y_x)_{x\in\Z^2}\) are iid. Each mark includes its entire future clock
histories.

The event \(A\) that the eventual configuration is all 1 is measurable
in this marked field: eventual states are limits of the measurable
integer-time states. It is invariant under all spatial translations,
since the graphical construction commutes with them.
We recall the elementary zero-one argument for this situation.
For any \(\epsilon>0\), choose an event \(B\) depending on finitely many
spatial marks with \(\Prob(A\mathbin{\triangle} B)<\epsilon\).
Such approximations hold for every event in a product sigma field:
the class of approximable events contains the finite-coordinate events
and is closed under complements and countable unions, the latter by
first taking a finite prefix with arbitrarily small probability error.
This also covers completed null events.

Choose a spatial translate \(\theta_zB\) whose coordinates are disjoint
from those of \(B\). The two events are independent, while
\(\theta_zA=A\). Consequently,
\[
 |\Prob(A)-\Prob(B)^2|\le2\epsilon,\qquad
 |\Prob(A)^2-\Prob(B)^2|\le2\epsilon.
\]
It follows that \(\Prob(A)(1-\Prob(A))\le4\epsilon\).
Letting \(\epsilon\) decrease to zero shows that \(\Prob(A)\) is 0 or 1.
Together with Lemma~\ref{lem:color}, this says that the limiting color
is deterministic.

If that color were 2, bounded convergence would give \(m(t)\to-1\).
But \eqref{eq:drift} gives \(m(t)\ge m(0)\) for every \(t\), and
\[
 m(0)=a_1-a_2=2a_1+a_0-1>-1.
\]
This contradiction proves that every site is eventually in state 1.
For finite nonempty \(K\), the maximum of its finitely many site
fixation times is almost surely finite. Its all-CO indicator therefore
tends to 1; bounded convergence gives \eqref{eq:target}. For empty \(K\)
the probability is identically 1. The bounded-convergence statements
apply along every real sequence tending to infinity and hence give
the full continuous-time limits.
\end{proof}

\appendix
\section{Graphical construction and the counting identity}\label{sec:clocks}

For completeness, we justify the two clock facts used in the proof.
The independent initial states and the exponential interarrival
variables of the countably many clocks can be constructed on a
countable product of real spaces. Their finite product laws are
consistent, so the countable extension theorem applies
\cite[Section~2.1.4, Theorem~2.1.21]{PTE}.
The exponential construction gives the required independent Poisson
histories \cite[Section~3.7]{PTE}.

Put \(\lambda=q/4+r/4+\gamma\), the total edge-clock rate.
For a fixed nearest-neighbor path of \(n\) steps and a fixed time
horizon \(T\), the expected number of its backward, strictly
time-ordered edge-ring sequences is
\(\lambda^nT^n/n!\).
This includes paths that repeat an edge. To verify the formula, first
count rings in distinct ordered dyadic time bins. Independent counts
in disjoint bins give the product of their means. Refinement increases
these counts to all strictly ordered tuples, so monotone convergence
gives the ordered time simplex of volume \(T^n/n!\).
There are \(4^n\) spatial paths from a site. Therefore
\[
 \Prob(\text{a backward edge-ring path of \(n\) steps before \(T\)})
       \le\frac{(4\lambda T)^n}{n!}\longrightarrow0.
\]

To determine a state, trace the most recent incident mark backwards
and, at an edge mark, inspect the two predecessor endpoint states.
The dependency tree branches at most twice. An infinite branch must
make infinitely many spatial steps, because each site has only finitely
many incident marks before \(T\). The preceding bound excludes such a
branch. A finitely branching infinite tree would contain an infinite
branch, so the dependency tree is finite almost surely.
Its leaves are initial states, and its updates uniquely determine the
state in question. Intersect over sites and integer horizons.
The resulting paths are right-continuous with left limits, their
changes occur only at incident marks, and they are measurable and
nonanticipating. The construction commutes with lattice translations.
No infinite-volume independence of the evolved states is required.

To prove \eqref{eq:counts}, partition \([0,T]\) dyadically and on
\((t_k,t_{k+1}]\) replace \(h(\eta_{t-}|_F)\) by
\(h(\eta_{t_k}|_F)\). This coefficient is measurable at time \(t_k\);
the future increment of the clock is independent of all past data and
has mean \(c(t_{k+1}-t_k)\). Thus the identity holds for each step
approximation. There are only finitely many marks incident to \(F\)
before \(T\), and almost surely none at a dyadic endpoint. The
approximations therefore converge to the left-state value at every
ring of \(C\) and at almost every time. If \(|h|\le M\), their clock
sums are bounded by \(M C(0,T]\), of expectation \(McT\), and their time
integrands are bounded by \(M\). Dominated convergence proves
\eqref{eq:counts}, for signed \(h\) as well. Values at fixed deterministic
times agree almost surely with their left limits.

\section*{AI use}
OpenAI Codex was used to develop and check the argument and to draft this
note. These checks are AI based; independent mathematical review is
requested.

\begin{thebibliography}{9}

\bibitem{DS}
R. Durrett and G. Swindle,
\emph{Coexistence results for catalysts},
Probability Theory and Related Fields \textbf{98} (1994), 489--515.
\href{https://doi.org/10.1007/BF01192836}{doi:10.1007/BF01192836}.

\bibitem{IPS}
R. Durrett,
\emph{Interacting Particle Systems: Ideas, Techniques, Applications},
manuscript, version PASTA0902, Section~5.4,
\url{https://sites.math.duke.edu/~rtd/PASTA/PASTA0902.pdf}
(accessed 15 October 2026).

\bibitem{PTE}
R. Durrett,
\emph{Probability: Theory and Examples},
fifth-edition manuscript, version 5, 11 January 2019,
\url{https://sites.math.duke.edu/~rtd/PTE/PTE5_011119.pdf}.

\end{thebibliography}
\end{document}
